% 图 79.2 —— 由 `checks/emit_hand.py` 自动拼出（probe 精测 + prims2 图元分解）
%%
% ⚠ 这是**稿子不是成品**：跑 `checks/checkhand.py 79.2` 看指标 + 看叠合图。
%%
%% ⚠ 待核：
%   · 本图 probe 报出阴影族 1 个 —— **本脚本不生成阴影**（要照原书手写）；prims2 的 strokes 里可能含阴影线，核对时留意别画重
%   · 可疑字形 1 项（空心圈 ⊙ 常被认成字母 o）—— 见文件末
%%
\documentclass[12pt]{standalone}
\usepackage{fontspec}
\setsansfont{Arial}
\newfontfamily\mfallback{DejaVu Sans}
\usepackage{tikz}
\begin{document}
\begin{tikzpicture}[x=1pt, y=-1pt, line width=5.0pt, line cap=round, line join=round]

  %% ════ 直线段（prims2 的 strokes）════
  \draw[line width=5.0pt] (744.0,290.0) -- (744.0,295.0);
  \draw[line width=6.0pt] (925.0,329.0) -- (933.0,324.0);
  \draw[line width=6.0pt] (540.0,307.0) -- (553.0,314.0);
  \draw[line width=8.7pt] (537.0,298.0) -- (539.0,306.0);
  \draw[line width=5.5pt] (785.8,476.0) -- (797.0,460.0);
  \draw[line width=3.4pt] (934.0,333.0) -- (935.0,330.0);
  \draw[line width=6.0pt] (786.0,457.0) -- (793.0,456.0);
  \draw[line width=5.3pt] (166.0,48.0) -- (165.0,53.0);
  \draw[line width=6.5pt] (154.0,52.0) -- (161.0,44.0);
  \draw[line width=7.0pt] (154.0,52.0) -- (164.0,54.0);
  \draw[line width=6.0pt] (734.0,296.0) -- (739.0,288.0);
  \draw[line width=6.0pt] (933.0,324.0) -- (936.0,319.0);
  \draw[line width=5.7pt] (933.0,324.0) -- (935.0,329.0);
  \draw[line width=5.5pt] (739.0,288.0) -- (741.0,281.0);
  \draw[line width=4.3pt] (739.0,288.0) -- (743.0,289.0);
  \draw[line width=5.0pt] (165.0,37.0) -- (162.0,43.0);
  \draw[line width=5.1pt] (797.0,459.0) -- (793.0,456.0);
  \draw[line width=5.7pt] (165.0,54.0) -- (170.0,56.0);
  \draw[line width=5.0pt] (537.0,307.0) -- (531.0,309.0);
  \draw[line width=6.0pt] (222.0,369.0) -- (368.0,313.0);
  \draw[line width=7.0pt] (368.0,313.0) -- (387.0,291.9);
  \draw[line width=6.0pt] (387.0,291.9) -- (394.0,269.7);
  \draw[line width=6.0pt] (394.0,269.7) -- (388.9,245.9);
  \draw[line width=6.0pt] (388.9,245.9) -- (376.1,232.1);
  \draw[line width=3.0pt] (215.0,196.0) -- (212.0,196.0);
  \draw[line width=6.8pt] (855.8,366.2) -- (862.0,364.0);
  \draw[line width=6.0pt] (969.4,347.4) -- (976.0,363.8);
  \draw[line width=6.0pt] (751.0,297.0) -- (745.0,296.0);
  \draw[line width=7.0pt] (933.0,333.0) -- (925.0,330.0);
  \draw[line width=6.0pt] (735.0,297.0) -- (744.0,296.0);
  \draw[line width=9.0pt] (869.0,350.0) -- (862.0,364.0);
  \draw[line width=5.1pt] (166.0,47.0) -- (162.0,44.0);
  \draw[line width=6.0pt] (218.0,204.0) -- (216.0,196.0);
  \draw[line width=4.3pt] (802.0,458.0) -- (798.0,459.0);
  \draw[line width=5.0pt] (938.0,338.0) -- (934.0,334.0);
  \draw[line width=3.0pt] (537.0,306.0) -- (487.0,280.0);
  \draw[line width=5.0pt] (441.0,291.0) -- (441.0,54.1);
  \draw[line width=4.0pt] (441.0,54.1) -- (438.9,45.9);
  \draw[line width=4.0pt] (438.9,45.9) -- (424.1,27.1);
  \draw[line width=4.0pt] (424.1,27.1) -- (402.8,18.0);
  \draw[line width=4.0pt] (402.8,18.0) -- (57.9,19.1);
  \draw[line width=4.5pt] (28.1,40.9) -- (23.0,53.8);
  \draw[line width=5.0pt] (23.0,53.8) -- (22.0,408.9);
  \draw[line width=5.0pt] (39.6,435.6) -- (60.5,443.0);
  \draw[line width=5.0pt] (60.5,443.0) -- (403.6,443.0);
  \draw[line width=4.0pt] (403.6,443.0) -- (424.3,432.7);
  \draw[line width=5.0pt] (424.3,432.7) -- (437.0,415.9);
  \draw[line width=4.0pt] (437.0,415.9) -- (441.0,400.5);
  \draw[line width=4.0pt] (441.0,400.5) -- (441.0,291.0);
  \draw[line width=4.0pt] (585.0,504.0) -- (586.0,209.2);
  \draw[line width=4.0pt] (586.0,209.2) -- (588.0,198.1);
  \draw[line width=4.0pt] (588.0,198.1) -- (602.8,175.2);
  \draw[line width=4.0pt] (602.8,175.2) -- (624.0,166.0);
  \draw[line width=5.0pt] (624.0,166.0) -- (988.0,167.0);
  \draw[line width=4.0pt] (988.0,167.0) -- (1013.0,184.0);
  \draw[line width=4.0pt] (1013.0,184.0) -- (1022.0,207.2);
  \draw[line width=4.0pt] (1022.0,207.2) -- (1021.0,567.9);
  \draw[line width=4.0pt] (1021.0,567.9) -- (1009.0,593.0);
  \draw[line width=5.0pt] (983.8,607.0) -- (627.9,607.9);
  \draw[line width=4.0pt] (627.9,607.9) -- (602.0,598.0);

  %% ════ 曲线（prims2 的 curves —— 三段式，坐标同为 (y,x)）════
  \draw[line width=6.0pt] (785.0,456.0) .. controls (781.3,436.0) and (762.5,419.8) .. (743.7,414.0);
  \draw[line width=6.0pt] (743.7,414.0) .. controls (719.7,407.0) and (702.6,427.0) .. (683.9,437.0);
  \draw[line width=6.0pt] (683.9,437.0) .. controls (666.7,444.1) and (645.1,438.8) .. (635.0,422.0);
  \draw[line width=6.0pt] (635.0,422.0) .. controls (611.9,359.0) and (689.8,322.9) .. (733.0,297.0);
  \draw[line width=6.0pt] (869.0,349.0) .. controls (869.6,321.7) and (866.3,288.7) .. (838.0,274.0);
  \draw[line width=6.0pt] (838.0,274.0) .. controls (806.6,262.2) and (773.2,276.0) .. (745.0,289.0);
  \draw[line width=6.0pt] (785.0,458.0) .. controls (783.2,464.0) and (771.2,475.9) .. (785.8,476.0);
  \draw[line width=6.0pt] (154.0,52.0) .. controls (113.9,63.4) and (43.6,84.8) .. (59.0,140.0);
  \draw[line width=6.0pt] (59.0,140.0) .. controls (83.2,188.7) and (141.3,150.0) .. (180.4,165.0);
  \draw[line width=6.0pt] (180.4,165.0) .. controls (193.9,169.9) and (207.4,180.5) .. (211.0,195.0);
  \draw[line width=6.0pt] (924.0,330.0) .. controls (904.0,326.1) and (886.8,340.9) .. (870.0,349.0);
  \draw[line width=7.0pt] (220.0,369.0) .. controls (214.5,367.1) and (207.8,372.5) .. (212.6,377.6);
  \draw[line width=7.0pt] (212.6,377.6) .. controls (218.1,380.1) and (222.9,375.8) .. (221.0,370.0);
  \draw[line width=8.0pt] (211.0,197.0) .. controls (203.9,202.2) and (212.3,212.7) .. (218.0,205.0);
  \draw[line width=6.0pt] (216.0,195.0) .. controls (239.7,158.7) and (273.2,122.3) .. (266.8,75.8);
  \draw[line width=6.0pt] (266.8,75.8) .. controls (253.9,32.9) and (200.9,37.3) .. (167.0,47.0);
  \draw[line width=6.0pt] (376.1,232.1) .. controls (328.6,207.4) and (274.2,201.8) .. (219.0,205.0);
  \draw[line width=6.0pt] (793.0,456.0) .. controls (809.1,423.7) and (831.3,394.2) .. (855.8,366.2);
  \draw[line width=6.0pt] (936.0,329.0) .. controls (948.9,331.0) and (960.7,336.6) .. (969.4,347.4);
  \draw[line width=6.0pt] (976.0,363.8) .. controls (974.9,452.7) and (866.1,445.2) .. (804.0,458.0);
  \draw[line width=6.0pt] (803.0,457.0) .. controls (824.6,427.7) and (857.8,400.6) .. (862.0,364.0);
  \draw[line width=5.0pt] (57.9,19.1) .. controls (46.0,22.2) and (34.3,29.6) .. (28.1,40.9);
  \draw[line width=4.0pt] (22.0,408.9) .. controls (25.0,419.1) and (31.1,428.2) .. (39.6,435.6);
  \draw[line width=4.0pt] (1009.0,593.0) .. controls (1001.7,599.4) and (993.7,605.6) .. (983.8,607.0);
  \draw[line width=4.0pt] (602.0,598.0) .. controls (576.3,574.4) and (587.3,536.3) .. (585.0,504.0);

  %% ════ 标签（probe 直接给的 \node 行）════
  %% ⚠ 全部补了 fill=white：文字在最上层，否则与曲线交叉干扰
     \node[anchor=center,inner sep=0pt,fill=white,font=\fontsize{29.7}{37.1}\selectfont] at (69.7,59.9) {{\mfallback\textbf{−}}};   % box(56,53,13,5)
     \node[anchor=center,inner sep=0pt,fill=white,font=\fontsize{30.7}{38.3}\selectfont] at (59.8,74.8) {\textsf{\textbf{\textit{b}}}};   % box(50,62,17,23)
     \node[anchor=center,inner sep=0pt,fill=white,font=\fontsize{31.4}{39.3}\selectfont] at (207.8,229.1) {\textsf{\textbf{\textit{e}}}};   % box(198,219,16,17)
     \node[anchor=center,inner sep=0pt,fill=white,font=\fontsize{21.5}{26.9}\selectfont] at (220.0,238.9) {\textsf{\textbf{1}}};   % box(214,231,7,15)
     \node[anchor=center,inner sep=0pt,fill=white,font=\fontsize{28.7}{35.9}\selectfont] at (410.5,269.0) {\textsf{\textbf{\textit{Y}}}};   % box(402,256,16,24)
     \node[anchor=center,inner sep=0pt,fill=white,font=\fontsize{29.2}{36.4}\selectfont] at (524.9,274.0) {\textsf{\textbf{\textit{P}}}};   % box(514,262,18,22)
     \node[anchor=center,inner sep=0pt,fill=white,font=\fontsize{32.6}{40.8}\selectfont] at (997.0,385.5) {\textsf{\textbf{\textit{α}}}};   % box(988,376,19,18)
     \node[anchor=center,inner sep=0pt,fill=white,font=\fontsize{30.0}{37.5}\selectfont] at (616.8,392.2) {\textsf{\textbf{\textit{b}}}};   % box(607,380,17,22)
     \node[anchor=center,inner sep=0pt,fill=white,font=\fontsize{32.0}{40.0}\selectfont] at (211.1,397.1) {\textsf{\textbf{\textit{a}}}};   % box(202,387,16,17)
     \node[anchor=center,inner sep=0pt,fill=white,font=\fontsize{24.1}{30.2}\selectfont] at (225.0,407.2) {\textsf{\textbf{0}}};   % box(217,399,12,16)
     \node[anchor=center,inner sep=0pt,fill=white,font=\fontsize{30.3}{37.9}\selectfont] at (448.7,434.0) {\textsf{\textbf{\textit{E}}}};   % box(439,422,20,22)
     \node[anchor=center,inner sep=0pt,fill=white,font=\fontsize{30.7}{38.3}\selectfont] at (780.8,500.8) {\textsf{\textbf{\textit{b}}}};   % box(771,488,17,23)
     \node[anchor=center,inner sep=0pt,fill=white,font=\fontsize{23.8}{29.8}\selectfont] at (795.5,512.7) {\textsf{\textbf{0}}};   % box(788,504,11,17)
     \node[anchor=center,inner sep=0pt,fill=white,font=\fontsize{30.3}{37.9}\selectfont] at (1038.8,595.5) {\textsf{\textbf{\textit{B}}}};   % box(1027,584,20,22)

  % 钉死画布：否则超出的图元/控制点会把页面撑大，1:1 回比就失准
  \pgfresetboundingbox
  \path[use as bounding box] (0,0) rectangle (1055,619);
\end{tikzpicture}
\end{document}

% ⚠ 待核清单（probe 拿不准，人眼过一遍）：
%   · 未识别/跳过：⑤ 跳过/未识别的字形
